\documentclass[11pt,a4paper]{article}
\usepackage[T1]{fontenc}
\usepackage{lmodern,array,amsmath,amssymb,booktabs,enumitem,fancyhdr,lastpage,needspace,microtype}
\usepackage[margin=19mm,headheight=14pt,footskip=11mm]{geometry}
\usepackage[hidelinks]{hyperref}
\setlength{\parindent}{0pt}
\setlength{\parskip}{5pt}
\setlength{\emergencystretch}{2em}
\pagestyle{fancy}\fancyhf{}
\fancyhead[L]{\small CS182 --- Lecture 04}
\fancyhead[R]{\small Solutions}
\fancyfoot[L]{\small Independent practice quiz}
\fancyfoot[R]{\small Page \thepage\ of \pageref{LastPage}}
\newcommand{\SetAnswer}[2]{\expandafter\gdef\csname QuizAnswer#1\endcsname{#2}}
\newcommand{\GetAnswer}[1]{\csname QuizAnswer#1\endcsname}
\newlength{\BlankAnswerHeight}\setlength{\BlankAnswerHeight}{16mm}
\newcommand{\AnswerBlank}{\leavevmode\par\vspace{\BlankAnswerHeight}}

\begin{document}
\begin{center}{\Large\bfseries CS182: Introduction to Machine Learning}\par{\Large\bfseries Lecture 04: Optimization Primer}\par{\large Worked Solutions and Marking Guide}\end{center}
\textbf{Scope:} Optimization\_Primer.pdf. Independent practice material based on the official 2026 Fall lecture.\par\textbf{Total:} 100 marks. \quad \textbf{Suggested time:} 120 minutes.
\begin{center}\begin{tabular}{lrrr}\toprule
Section & Questions & Each & Marks\\\midrule
A. Misconceptions: True/False & 12 & 1 & 12\\
B. Multiple response & 12 & 2 & 24\\
C. Proofs and formula derivations & 4 & 10 & 40\\
D. Applications & 2 & 12 & 24\\\midrule
Total & 30 & & 100\\\bottomrule\end{tabular}\end{center}
\textbf{Scoring.} In Section B each item has exactly four options, of which one,
two, or three are correct. The complete correct set earns 2 marks; a nonempty
proper subset earns 1 mark; any incorrect option or no selection earns 0.
T/F answers earn 1 mark only when correct. Show reasoning in Sections C and D;
part marks are available for valid intermediate work. Equivalent correct methods receive credit.
\par\medskip\textbf{Instructions.} Answer in English. Use natural logarithms.
Exact values and justified approximations are accepted. All vectors are columns
unless stated otherwise. Use local assumptions, label conventions, and tie rules.
\par\medskip\textbf{Quick key --- Section A}
\begin{center}\renewcommand{\arraystretch}{1.15}\begin{tabular}{|c|c|c|c|c|c|}
\hline
1 & 2 & 3 & 4 & 5 & 6\\\hline
\makebox[20mm]{\rule{0pt}{5mm}T} & \makebox[20mm]{\rule{0pt}{5mm}T} & \makebox[20mm]{\rule{0pt}{5mm}T} & \makebox[20mm]{\rule{0pt}{5mm}F} & \makebox[20mm]{\rule{0pt}{5mm}T} & \makebox[20mm]{\rule{0pt}{5mm}F}\\\hline
\end{tabular}\end{center}
\begin{center}\renewcommand{\arraystretch}{1.15}\begin{tabular}{|c|c|c|c|c|c|}
\hline
7 & 8 & 9 & 10 & 11 & 12\\\hline
\makebox[20mm]{\rule{0pt}{5mm}T} & \makebox[20mm]{\rule{0pt}{5mm}F} & \makebox[20mm]{\rule{0pt}{5mm}T} & \makebox[20mm]{\rule{0pt}{5mm}F} & \makebox[20mm]{\rule{0pt}{5mm}T} & \makebox[20mm]{\rule{0pt}{5mm}F}\\\hline
\end{tabular}\end{center}
\par\medskip\textbf{Quick key --- Section B}
\begin{center}\renewcommand{\arraystretch}{1.15}\begin{tabular}{|c|c|c|c|c|c|}
\hline
13 & 14 & 15 & 16 & 17 & 18\\\hline
\makebox[20mm]{\rule{0pt}{5mm}BC} & \makebox[20mm]{\rule{0pt}{5mm}AC} & \makebox[20mm]{\rule{0pt}{5mm}ABC} & \makebox[20mm]{\rule{0pt}{5mm}CD} & \makebox[20mm]{\rule{0pt}{5mm}AD} & \makebox[20mm]{\rule{0pt}{5mm}AC}\\\hline
\end{tabular}\end{center}
\begin{center}\renewcommand{\arraystretch}{1.15}\begin{tabular}{|c|c|c|c|c|c|}
\hline
19 & 20 & 21 & 22 & 23 & 24\\\hline
\makebox[20mm]{\rule{0pt}{5mm}BD} & \makebox[20mm]{\rule{0pt}{5mm}BC} & \makebox[20mm]{\rule{0pt}{5mm}BCD} & \makebox[20mm]{\rule{0pt}{5mm}A} & \makebox[20mm]{\rule{0pt}{5mm}B} & \makebox[20mm]{\rule{0pt}{5mm}CD}\\\hline
\end{tabular}\end{center}
\newpage\section*{Section A. True/False --- Explanations}
\par\noindent\begin{minipage}{\linewidth}
\textbf{1.} \textbf{True.} Each convex sublevel set is convex, affine equality sets are convex, and intersections preserve convexity.

{\footnotesize\textit{Source: PDF pp.16-18.}}
\end{minipage}\par\vspace{3mm}
\par\noindent\begin{minipage}{\linewidth}
\textbf{2.} \textbf{True.} The Hessian is the constant matrix \(Q\); apply the second-order convexity criterion on \(\mathbb R^n\).

{\footnotesize\textit{Source: PDF pp.18,20.}}
\end{minipage}\par\vspace{3mm}
\par\noindent\begin{minipage}{\linewidth}
\textbf{3.} \textbf{True.} For a fixed primal variable, the Lagrangian is affine in the multipliers. Taking its pointwise infimum produces a concave function, possibly with value \(-\infty\).

{\footnotesize\textit{Source: PDF pp.5,21.}}
\end{minipage}\par\vspace{3mm}
\par\noindent\begin{minipage}{\linewidth}
\textbf{4.} \textbf{False.} Complementary slackness only states \(\lambda_if_i(x)=0\). For minimizing \(x^2\) subject to \(-x\le0\), \(x=0\) is active and the KKT multiplier is zero.

{\footnotesize\textit{Source: PDF pp.3,11-12.}}
\end{minipage}\par\vspace{3mm}
\par\noindent\begin{minipage}{\linewidth}
\textbf{5.} \textbf{True.} Weak duality gives \(g(\lambda,\nu)\le d^*\le p^*\le f_0(x)\). Equality of the endpoints makes every inequality an equality.

{\footnotesize\textit{Source: PDF pp.6-10,13-14.}}
\end{minipage}\par\vspace{3mm}
\par\noindent\begin{minipage}{\linewidth}
\textbf{6.} \textbf{False.} For \(f(x)=-x^2\), the gradient vanishes at zero, which is a global maximum. Convexity requires the global convexity inequality; for twice continuously differentiable functions on an open convex domain, it is equivalent to a positive semidefinite Hessian throughout that domain.

{\footnotesize\textit{Source: PDF pp.18,22-24.}}
\end{minipage}\par\vspace{3mm}
\par\noindent\begin{minipage}{\linewidth}
\textbf{7.} \textbf{True.} If a feasible point had a lower objective, points on the segment toward it would be feasible, arbitrarily close to the local minimizer, and have a lower objective by convexity.

{\footnotesize\textit{Source: PDF pp.16-17,22.}}
\end{minipage}\par\vspace{3mm}
\par\noindent\begin{minipage}{\linewidth}
\textbf{8.} \textbf{False.} For \(f(x)=x^2/2\), the update is \(x_{k+1}=(1-\eta)x_k\). Taking \(\eta=3\) makes nonzero iterates grow in magnitude. A step-size restriction is needed.

{\footnotesize\textit{Source: PDF p.24.}}
\end{minipage}\par\vspace{3mm}
\par\noindent\begin{minipage}{\linewidth}
\textbf{9.} \textbf{True.} If \(h_j(x)=0\), then \(\nu_jh_j(x)=0\) for every real \(\nu_j\). Inequality multipliers have a different sign restriction.

{\footnotesize\textit{Source: PDF pp.13-14.}}
\end{minipage}\par\vspace{3mm}
\par\noindent\begin{minipage}{\linewidth}
\textbf{10.} \textbf{False.} Consider minimizing \(x\) subject to \(x^2\le0\). The sole feasible point and optimum is zero, but stationarity would require \(1+2\lambda\cdot0=0\), which no finite multiplier satisfies. Regularity conditions matter when asserting multiplier existence or dual attainment.

{\footnotesize\textit{Source: PDF pp.8,12,14-18.}}
\end{minipage}\par\vspace{3mm}
\par\noindent\begin{minipage}{\linewidth}
\textbf{11.} \textbf{True.} With \(\lambda_i>0\), complementary slackness forces \(f_i(x)=0\).

{\footnotesize\textit{Source: PDF pp.11-12,15.}}
\end{minipage}\par\vspace{3mm}
\par\noindent\begin{minipage}{\linewidth}
\textbf{12.} \textbf{False.} The lower-bound proof uses \(\lambda_i\ge0\), ensuring \(\lambda_if_i(x)\le0\) at a feasible point. Negative inequality multipliers generally break that step.

{\footnotesize\textit{Source: PDF pp.4-8.}}
\end{minipage}\par\vspace{3mm}
\newpage\section*{Section B. Multiple Response --- Explanations}
\par\noindent\begin{minipage}{\linewidth}
% BEGIN QUESTION Q13 / 题目与选项区
\textbf{13.} \textbf{Correct set: B, C.}
\begin{enumerate}[label=(\Alph*),leftmargin=8mm,itemsep=2pt,topsep=2pt]
\item \textbf{Incorrect.} Convexity requires positive semidefiniteness. A linear function has a zero Hessian.
\item \textbf{Correct.} The first-order inequality reduces to \(f(y)\ge f(x)\) for every \(y\).
\item \textbf{Correct.} The tangent affine function is a global lower bound by the first-order criterion.
\item \textbf{Incorrect.} A constant function is convex and has infinitely many global minimizers.
\end{enumerate}
{\footnotesize\textit{Source: PDF pp.18,22.}}
% END QUESTION Q13

\end{minipage}\par\vspace{3mm}
\par\noindent\begin{minipage}{\linewidth}
% BEGIN QUESTION Q14 / 题目与选项区
\textbf{14.} \textbf{Correct set: A, C.}
\begin{enumerate}[label=(\Alph*),leftmargin=8mm,itemsep=2pt,topsep=2pt]
\item \textbf{Correct.} It is a pointwise infimum of affine functions of the multipliers.
\item \textbf{Incorrect.} For example, \((x^2-1)^2\) on the unconstrained real line has derivative zero at zero, yet lower values at \(\pm1\).
\item \textbf{Correct.} Weak duality follows directly from feasibility and the infimum definition.
\item \textbf{Incorrect.} Without convexity, stationarity can describe a maximum or a saddle point.
\end{enumerate}
{\footnotesize\textit{Source: PDF pp.6-8,12,21-24.}}
% END QUESTION Q14

\end{minipage}\par\vspace{3mm}
\par\noindent\begin{minipage}{\linewidth}
% BEGIN QUESTION Q15 / 题目与选项区
\textbf{15.} \textbf{Correct set: A, B, C.}
\begin{enumerate}[label=(\Alph*),leftmargin=8mm,itemsep=2pt,topsep=2pt]
\item \textbf{Correct.} Substituting \(x=3-\lambda/2\) gives \(\lambda^2/4+\lambda(2-\lambda/2)\).
\item \textbf{Correct.} Stationarity gives \(2(x-3)+\lambda=0\); the second derivative is positive.
\item \textbf{Correct.} The concave quadratic has derivative \(2-\lambda/2\), so its maximum on \(\lambda\ge0\) occurs at four.
\item \textbf{Incorrect.} It violates the nonnegative multiplier restriction for the inequality, even though its infimum is finite.
\end{enumerate}
{\footnotesize\textit{Source: PDF pp.4-8,11-12.}}
% END QUESTION Q15

\end{minipage}\par\vspace{3mm}
\par\noindent\begin{minipage}{\linewidth}
% BEGIN QUESTION Q16 / 题目与选项区
\textbf{16.} \textbf{Correct set: C, D.}
\begin{enumerate}[label=(\Alph*),leftmargin=8mm,itemsep=2pt,topsep=2pt]
\item \textbf{Incorrect.} The feasible set is convex, but the minimization objective is concave and not convex.
\item \textbf{Incorrect.} The equality produces the nonconvex feasible set \(\{-1,1\}\); nonlinear equalities cannot generally be treated as affine constraints.
\item \textbf{Correct.} The objective is convex, the equality is affine, and the inequality defines a half-space.
\item \textbf{Correct.} The objective is affine and the inequality is a sublevel set of the convex function \(x^2+y^2-4\).
\end{enumerate}
{\footnotesize\textit{Source: PDF pp.16-20.}}
% END QUESTION Q16

\end{minipage}\par\vspace{3mm}
\newpage\section*{Section B. Multiple Response --- Explanations}
\par\noindent\begin{minipage}{\linewidth}
% BEGIN QUESTION Q17 / 题目与选项区
\textbf{17.} \textbf{Correct set: A, D.}
\begin{enumerate}[label=(\Alph*),leftmargin=8mm,itemsep=2pt,topsep=2pt]
\item \textbf{Correct.} The scalar quadratic uses the symmetric part \((G+G^T)/2\), giving the displayed derivative.
\item \textbf{Incorrect.} The symmetric Hessian has determinant \(8-9=-1\), so it is indefinite.
\item \textbf{Incorrect.} The Hessian is \((G+G^T)/2=\begin{pmatrix}2&3\\3&4\end{pmatrix}\).
\item \textbf{Correct.} The indefinite symmetric Hessian has both positive and negative eigenvalues, so the quadratic takes both signs arbitrarily close to zero.
\end{enumerate}
{\footnotesize\textit{Source: PDF pp.18,20,24.}}
% END QUESTION Q17

\end{minipage}\par\vspace{3mm}
\par\noindent\begin{minipage}{\linewidth}
% BEGIN QUESTION Q18 / 题目与选项区
\textbf{18.} \textbf{Correct set: A, C.}
\begin{enumerate}[label=(\Alph*),leftmargin=8mm,itemsep=2pt,topsep=2pt]
\item \textbf{Correct.} A norm measures how close the gradient vector is to zero; its interpretation as an objective-error certificate requires extra assumptions.
\item \textbf{Incorrect.} For \(f(x)=a x^2/2\), fixing a small nonzero gradient \(g=ax\) gives objective error \(g^2/(2a)\), which can be arbitrarily large as \(a\) decreases.
\item \textbf{Correct.} The first-order lower bound gives \(f(y)\ge f(x)\) for every \(y\).
\item \textbf{Incorrect.} The constraint gradients can balance the objective gradient through the KKT stationarity equation.
\end{enumerate}
{\footnotesize\textit{Source: PDF pp.12,15,18,24.}}
% END QUESTION Q18

\end{minipage}\par\vspace{3mm}
\par\noindent\begin{minipage}{\linewidth}
% BEGIN QUESTION Q19 / 题目与选项区
\textbf{19.} \textbf{Correct set: B, D.}
\begin{enumerate}[label=(\Alph*),leftmargin=8mm,itemsep=2pt,topsep=2pt]
\item \textbf{Incorrect.} The dual function takes the infimum over the primal domain, without retaining the constraints that were moved into the Lagrangian.
\item \textbf{Correct.} It seeks the best lower bound on the primal minimization optimum.
\item \textbf{Incorrect.} The lecture defines dual feasibility to require \(g> -\infty\) in addition to nonnegative inequality multipliers.
\item \textbf{Correct.} This sign convention makes each inequality term nonpositive at a feasible point.
\end{enumerate}
{\footnotesize\textit{Source: PDF pp.4-8,13-14.}}
% END QUESTION Q19

\end{minipage}\par\vspace{3mm}
\par\noindent\begin{minipage}{\linewidth}
% BEGIN QUESTION Q20 / 题目与选项区
\textbf{20.} \textbf{Correct set: B, C.}
\begin{enumerate}[label=(\Alph*),leftmargin=8mm,itemsep=2pt,topsep=2pt]
\item \textbf{Incorrect.} There are \(m\) multipliers, one for each primal inequality, and \(n\) equality constraints.
\item \textbf{Correct.} The Lagrangian is \((c+A^T\lambda)^Tx-b^T\lambda\); the unrestricted infimum is finite only when the coefficient of \(x\) vanishes.
\item \textbf{Correct.} Move \(x\) arbitrarily far in the negative direction of its nonzero linear coefficient.
\item \textbf{Incorrect.} The sign of inequality multipliers follows from the primal inequality direction, independently of the domain of \(x\).
\end{enumerate}
{\footnotesize\textit{Source: PDF pp.5,9.}}
% END QUESTION Q20

\end{minipage}\par\vspace{3mm}
\newpage\section*{Section B. Multiple Response --- Explanations}
\par\noindent\begin{minipage}{\linewidth}
% BEGIN QUESTION Q21 / 题目与选项区
\textbf{21.} \textbf{Correct set: B, C, D.}
\begin{enumerate}[label=(\Alph*),leftmargin=8mm,itemsep=2pt,topsep=2pt]
\item \textbf{Incorrect.} The second-coordinate factor has magnitude \(1.4>1\), so any nonzero second coordinate diverges.
\item \textbf{Correct.} The gradient is \((x_1,4x_2)^T\); subtracting \(0.6\) times it gives the factors.
\item \textbf{Correct.} The second-coordinate multiplier is the negative number \(-1.4\).
\item \textbf{Correct.} The second coordinate stays zero and the first is multiplied repeatedly by \(0.4\).
\end{enumerate}
{\footnotesize\textit{Source: PDF pp.18,20,24.}}
% END QUESTION Q21

\end{minipage}\par\vspace{3mm}
\par\noindent\begin{minipage}{\linewidth}
% BEGIN QUESTION Q22 / 题目与选项区
\textbf{22.} \textbf{Correct set: A.}
\begin{enumerate}[label=(\Alph*),leftmargin=8mm,itemsep=2pt,topsep=2pt]
\item \textbf{Correct.} It is a special case of the pointwise supremum rule.
\item \textbf{Incorrect.} The minimum of the affine functions \(x\) and \(-x\) is \(-|x|\), which is not convex.
\item \textbf{Incorrect.} Only the weighted convexity inequality and existence of the expectations are needed.
\item \textbf{Incorrect.} Jensen gives the reverse inequality. For \(f(x)=x^2\) and a fair sign \(X\in\{-1,1\}\), the displayed claim reads \(0\ge1\), which is false.
\end{enumerate}
{\footnotesize\textit{Source: PDF pp.17-18.}}
% END QUESTION Q22

\end{minipage}\par\vspace{3mm}
\par\noindent\begin{minipage}{\linewidth}
% BEGIN QUESTION Q23 / 题目与选项区
\textbf{23.} \textbf{Correct set: B.}
\begin{enumerate}[label=(\Alph*),leftmargin=8mm,itemsep=2pt,topsep=2pt]
\item \textbf{Incorrect.} The point is feasible and complementary slackness holds, but stationarity fails: \(2(0)-1-2(0)=-1\ne0\).
\item \textbf{Correct.} Both constraint values equal zero at the point.
\item \textbf{Incorrect.} The second constraint has gradient \(-2\), so the correct equation is \(2x-\lambda_1-2\lambda_2=0\).
\item \textbf{Incorrect.} Stationarity requires \(-\lambda_1-2\lambda_2=0\); nonnegativity then forces both multipliers to be zero.
\end{enumerate}
{\footnotesize\textit{Source: PDF pp.3,11-12.}}
% END QUESTION Q23

\end{minipage}\par\vspace{3mm}
\par\noindent\begin{minipage}{\linewidth}
% BEGIN QUESTION Q24 / 题目与选项区
\textbf{24.} \textbf{Correct set: C, D.}
\begin{enumerate}[label=(\Alph*),leftmargin=8mm,itemsep=2pt,topsep=2pt]
\item \textbf{Incorrect.} The interval bounds the optimum but does not identify its midpoint as optimal.
\item \textbf{Incorrect.} The gap bounds objective error. A distance guarantee needs additional curvature information and a corresponding conversion.
\item \textbf{Correct.} From \(p^*\ge10.9\), \(11.2-p^*\le0.3\).
\item \textbf{Correct.} The dual value is a lower bound and the feasible primal value is an upper bound.
\end{enumerate}
{\footnotesize\textit{Source: PDF pp.6-10.}}
% END QUESTION Q24

\end{minipage}\par\vspace{3mm}
\newpage\section*{Section C. Proofs and Formula Derivations}
\textbf{25. Proving the first-order characterization of convexity}\hfill\textbf{[10 marks]}\par\medskip
\textbf{(a) [5 marks]} For \(0<t\le1\), convexity gives \(f(x+t(y-x))\le(1-t)f(x)+tf(y)\). Rearranging yields \[\frac{f(x+t(y-x))-f(x)}{t}\le f(y)-f(x).\] Differentiability makes the left side converge to \(\nabla f(x)^T(y-x)\) as \(t\downarrow0\). Hence \(f(y)\ge f(x)+\nabla f(x)^T(y-x)\). Convexity of the domain ensures the entire segment remains in the domain.

\textbf{(b) [5 marks]} Applying the assumed inequality at \(z\) gives \(f(x)\ge f(z)+\nabla f(z)^T(x-z)\) and \(f(y)\ge f(z)+\nabla f(z)^T(y-z)\). Multiply these by \(1-t\) and \(t\), respectively, and add. Since \((1-t)(x-z)+t(y-z)=0\), the gradient terms cancel: \[(1-t)f(x)+tf(y)\ge f(z)=f((1-t)x+ty).\] At \(t=0\) or \(t=1\), equality is immediate. This proves convexity.

\textit{Marking: use the displayed subpart allocations; award credit for the logical or numerical intermediate steps.}\par
{\footnotesize\textit{Source: PDF pp.17-18.}}
\newpage\section*{Section C. Proofs and Formula Derivations}
\textbf{26. Dual concavity, weak duality, and complementary slackness}\hfill\textbf{[10 marks]}\par\medskip
\textbf{(a) [4 marks]} Write a combined multiplier vector as \(u=(\lambda,\nu)\). For \(0\le t\le1\), affinity in multipliers gives \(L(x,tu+(1-t)v)=tL(x,u)+(1-t)L(x,v)\). For every \(x\), the right side is at least \(t\inf_zL(z,u)+(1-t)\inf_zL(z,v)=tg(u)+(1-t)g(v)\). Taking the infimum over \(x\) yields \[g(tu+(1-t)v)\ge tg(u)+(1-t)g(v).\] No convexity of the functions in \(x\) was used.

\textbf{(b) [3 marks]} Primal feasibility gives \(f_i(x)\le0\) and \(h_j(x)=0\). Thus \[g(\lambda,\nu)\le L(x,\lambda,\nu)=f_0(x)+\sum_i\lambda_if_i(x)\le f_0(x).\] The last inequality uses \(\lambda_i\ge0\); equality terms vanish for every real \(\nu_j\). Taking the infimum over feasible \(x\) and the supremum over dual feasible multipliers gives \(d^*\le p^*\).

\textbf{(c) [3 marks]} The inequalities from the preceding proof become \[g(\lambda^*,\nu^*)\le L(x^*,\lambda^*,\nu^*)\le f_0(x^*)=g(\lambda^*,\nu^*).\] Therefore both inequalities are equalities. In particular, \(\sum_i\lambda_i^*f_i(x^*)=0\). Each summand is nonpositive, so each must be zero: \(\lambda_i^*f_i(x^*)=0\). The other equality, \(L(x^*,\lambda^*,\nu^*)=\inf_xL(x,\lambda^*,\nu^*)\), says precisely that \(x^*\) globally minimizes the Lagrangian.

\textit{Marking: use the displayed subpart allocations; award credit for the logical or numerical intermediate steps.}\par
{\footnotesize\textit{Source: PDF pp.4-8,11,13-14,21.}}
\newpage\section*{Section C. Proofs and Formula Derivations}
\textbf{27. Why convex KKT conditions are sufficient}\hfill\textbf{[10 marks]}\par\medskip
\textbf{(a) [4 marks]} Because \(\lambda_i^*\ge0\), the weighted sum \(f_0+\sum_i\lambda_i^*f_i\) is convex. Each equality contribution \(\nu_j^*h_j\) is affine even when \(\nu_j^*\) is negative. Thus \(L(\cdot,\lambda^*,\nu^*)\) is differentiable convex. Stationarity states that its gradient at \(x^*\) is zero. Applying the first-order inequality gives \[L(x,\lambda^*,\nu^*)\ge L(x^*,\lambda^*,\nu^*)+0^T(x-x^*)\] for every \(x\), proving global minimality.

\textbf{(b) [3 marks]} Feasibility and complementary slackness imply \(L(x^*,\lambda^*,\nu^*)=f_0(x^*)\). Since the previous part establishes a global minimum, \(g(\lambda^*,\nu^*)=L(x^*,\lambda^*,\nu^*)=f_0(x^*)\). The finite equality of a dual lower bound and a feasible primal upper bound proves both optimalities and zero gap by weak duality. No constraint qualification is needed for this sufficiency implication; such conditions are used to establish that suitable multipliers exist at an optimum.

\textbf{(c) [3 marks]} Convex inequalities and affine equalities define a convex feasible set. Suppose distinct feasible minimizers \(u,v\) both have value \(p^*\). Their midpoint \(w=(u+v)/2\) is feasible. Strict convexity gives \[f_0(w)<\tfrac12f_0(u)+\tfrac12f_0(v)=p^*,\] contradicting optimality. The convex feasible set is needed to ensure the midpoint is an admissible competitor.

\textit{Marking: use the displayed subpart allocations; award credit for the logical or numerical intermediate steps.}\par
{\footnotesize\textit{Source: PDF pp.12-18,22; sufficiency and uniqueness are derived from the lecture conditions.}}
\newpage\section*{Section C. Proofs and Formula Derivations}
\textbf{28. Deriving a quadratic gradient-descent step-size bound}\hfill\textbf{[10 marks]}\par\medskip
\textbf{(a) [3 marks]} Symmetry gives \[f(x+h)=f(x)+(Qx+c)^Th+\tfrac12h^TQh.\] Therefore \(\nabla f(x)=Qx+c\) and \(\nabla^2f(x)=Q\). The second-order criterion gives convexity exactly when \(Q\succeq0\).

\textbf{(b) [4 marks]} Substitute \(h=-\eta g\) in the exact expansion: \[f(x^+)-f(x)=-\eta\|g\|^2+\tfrac12\eta^2g^TQg.\] Since \(Q\preceq LI\), \(g^TQg\le L\|g\|^2\), so \[f(x^+)-f(x)\le-\eta\left(1-\frac{\eta L}{2}\right)\|g\|^2.\] For \(0<\eta<2/L\), the coefficient is strictly negative whenever the gradient is nonzero. At \(\eta=2/L\), this estimate guarantees only nonincrease, not strict decrease or convergence.

\textbf{(c) [3 marks]} The unique stationary point is \(x^*=-Q^{-1}c\). The error satisfies \(e_{k+1}=(I-\eta Q)e_k\). Diagonalize \(Q=U\operatorname{diag}(\lambda_1,\ldots,\lambda_n)U^T\), with all \(\lambda_i>0\). In coordinates \(z_k=U^Te_k\), each error coordinate is multiplied by \(1-\eta\lambda_i\). The stated interval ensures \(|1-\eta\lambda_i|<1\) for every coordinate, so \(z_k\to0\) and \(x_k\to x^*\). At the endpoint for \(f(x)=Lx^2/2\), the factor is \(-1\), so a nonzero start alternates forever.

\textit{Marking: use the displayed subpart allocations; award credit for the logical or numerical intermediate steps.}\par
{\footnotesize\textit{Source: PDF pp.18,20,24; the descent bound and spectral convergence argument are supplemental derivations.}}
\newpage\section*{Section D. Applications}
\textbf{29. Capacity planning with an active resource cap}\hfill\textbf{[12 marks]}\par\medskip
\textbf{(a) [4 marks]} The equality gives \(y=4-x\). Nonnegativity and the cap give \(0\le x\le2.5\). The restricted objective is \((x-5)^2+2(2-x)^2=3x^2-18x+33\), with derivative \(6x-18\). Its unconstrained minimizer is \(x=3\), outside the interval; the function is decreasing throughout the interval. Therefore \(x^*=2.5\), \(y^*=1.5\), and \(f^*=(-2.5)^2+2(-0.5)^2=6.75=27/4\). The objective Hessian is \(\operatorname{diag}(2,4)\succ0\), so this feasible minimizer is unique.

\textbf{(b) [4 marks]} The Lagrangian is \[L=(x-5)^2+2(y-2)^2-\alpha x-\beta y+\lambda(x-2.5)+\nu(x+y-4).\] Since \(x^*,y^*>0\), complementary slackness gives \(\alpha=\beta=0\). The stationarity equations at \((2.5,1.5)\) are \(-5+\lambda+\nu=0\) and \(-2+\nu=0\); hence \(\nu=2\), \(\lambda=3\). These multipliers have the required signs, and the cap is active. For these multipliers, minimizing the strictly convex Lagrangian over unrestricted \(x,y\) gives the same \((2.5,1.5)\). The constraint contributions there vanish, so \(g(0,0,3,2)=6.75=f(x^*,y^*)\), certifying optimality.

\textbf{(c) [4 marks]} The primal objective is \(f(2,2)=9\). With the reported multipliers, Lagrangian stationarity gives \(2(x-5)+2+2=0\) and \(4(y-2)+2=0\), so its unconstrained minimum occurs at \((3,1.5)\). This point need not satisfy the capacity cap when evaluating the dual function. Its objective is \(4.5\); the multiplier terms are \(2(3-2.5)+2(3+1.5-4)=1+1=2\). Thus \(g=6.5\), and weak duality certifies \(p^*\in[6.5,9]\). The reported plan is at most \(9-6.5=2.5\) above optimum. Its actual error, using the optimum found earlier, is \(9-6.75=2.25\), consistent with the bound.

\textit{Marking: use the displayed subpart allocations; award credit for the logical or numerical intermediate steps.}\par
{\footnotesize\textit{Source: PDF pp.2-15,17-20.}}
\newpage\section*{Section D. Applications}
\textbf{30. Choosing a stable learning rate for a quadratic calibration loss}\hfill\textbf{[12 marks]}\par\medskip
\textbf{(a) [4 marks]} The gradient is \((2u-4,8v+8)^T\) and the Hessian is \(Q=\operatorname{diag}(2,8)\succ0\). Solving the gradient equations gives \((u^*,v^*)=(2,-1)\). Substitution gives \(f^*=\tfrac12(8+8)-8-8=-8\). Positive definiteness makes the objective strictly convex, so the stationary point is the unique global minimizer.

\textbf{(b) [4 marks]} The initial gradient is \((-4,8)^T\), so \((u_1,v_1)=(4\eta,-8\eta)\), and \(f(0,0)=0\). For \(\eta=0.2\), the point is \((0.8,-1.6)\) and the objective is \(0.64+10.24-3.2-12.8=-5.12\), an improvement with objective error \(2.88\). For \(\eta=0.3\), the point is \((1.2,-2.4)\) and the objective is \(1.44+23.04-4.8-19.2=0.48\), an increase with error \(8.48\). Both comparisons use the same exact quadratic loss.

\textbf{(c) [4 marks]} The error updates are \(u_{k+1}-2=(1-2\eta)(u_k-2)\) and \(v_{k+1}+1=(1-8\eta)(v_k+1)\). Both contraction conditions require \(0<\eta<2/8=0.25\). At \(\eta=0.25\), \(u_k=2-2(1/2)^k\) tends to two, while \(v_k=-1+(-1)^k\) alternates between zero and minus two. Completing squares gives \(f(u,v)=-8+(u-2)^2+4(v+1)^2\), so the objective tends to \(-4\), with successive changes tending to zero, despite failure of the iterates to converge or reach the optimum. A small objective change alone can therefore miss an oscillating, nonoptimal state. For an exact constant-objective counterexample, start instead at \((u_0,v_0)=(2,0)\). Then \(u_k=2\) and \(v_k\) alternates between zero and \(-2\), so \(f(u_k,v_k)=-4\) at every step even though \(f^*=-8\).

\textit{Marking: use the displayed subpart allocations; award credit for the logical or numerical intermediate steps.}\par
{\footnotesize\textit{Source: PDF pp.18,20,24; stability calculations apply the supplemental quadratic bound.}}
\newpage\section*{Lecture Alignment and Coverage}
Primary source: \texttt{Optimization\_Primer.pdf} in \texttt{E:\textbackslash CS182\textbackslash 2026fall\textbackslash lecture}. All PDF pages count the cover as page 1. Additional arguments are marked as supplementary derivations when appropriate. These are original practice questions, not copied homework questions.

\textbf{T/F}
\begin{description}[leftmargin=12mm,itemsep=1pt]
\item[1] PDF pp.16-18.
\item[2] PDF pp.18,20.
\item[3] PDF pp.5,21.
\item[4] PDF pp.3,11-12.
\item[5] PDF pp.6-10,13-14.
\item[6] PDF pp.18,22-24.
\item[7] PDF pp.16-17,22.
\item[8] PDF p.24.
\item[9] PDF pp.13-14.
\item[10] PDF pp.8,12,14-18.
\item[11] PDF pp.11-12,15.
\item[12] PDF pp.4-8.
\end{description}
\textbf{Multiple response}
\begin{description}[leftmargin=12mm,itemsep=1pt]
\item[13] PDF pp.18,22.
\item[14] PDF pp.6-8,12,21-24.
\item[15] PDF pp.4-8,11-12.
\item[16] PDF pp.16-20.
\item[17] PDF pp.18,20,24.
\item[18] PDF pp.12,15,18,24.
\item[19] PDF pp.4-8,13-14.
\item[20] PDF pp.5,9.
\item[21] PDF pp.18,20,24.
\item[22] PDF pp.17-18.
\item[23] PDF pp.3,11-12.
\item[24] PDF pp.6-10.
\end{description}
\textbf{Proofs and applications}
\begin{description}[leftmargin=12mm,itemsep=1pt]
\item[25] PDF pp.17-18.
\item[26] PDF pp.4-8,11,13-14,21.
\item[27] PDF pp.12-18,22; sufficiency and uniqueness are derived from the lecture conditions.
\item[28] PDF pp.18,20,24; the descent bound and spectral convergence argument are supplemental derivations.
\item[29] PDF pp.2-15,17-20.
\item[30] PDF pp.18,20,24; stability calculations apply the supplemental quadratic bound.
\end{description}
\end{document}
